%; whizzy paragraph -pdf xpdf -latex ./whizzypdfptex.sh
%; whizzy-paragraph "^\\\\begin{frame}\\|\\\\emtext"
% latex beamer presentation.
% platex, latex-beamer でコンパイルすることを想定。 

%     Tokyo Debian Meeting resources
%     Copyright (C) 2012 Junichi Uekawa

%     This program is free software; you can redistribute it and/or modify
%     it under the terms of the GNU General Public License as published by
%     the Free Software Foundation; either version 2 of the License, or
%     (at your option) any later version.

%     This program is distributed in the hope that it will be useful,
%     but WITHOUT ANY WARRANTY; without even the implied warreanty of
%     MERCHANTABILITY or FITNESS FOR A PARTICULAR PURPOSE.  See the
%     GNU General Public License for more details.
%     You should have received a copy of the GNU General Public License
%     along with this program; if not, write to the Free Software
%     Foundation, Inc., 51 Franklin St, Fifth Floor, Boston, MA  02110-1301 USA

\documentclass[cjk,dvipdfmx,12pt]{beamer}
\usetheme{Tokyo}
\usepackage{monthlypresentation}

%  preview (shell-command (concat "evince " (replace-regexp-in-string "tex$" "pdf"(buffer-file-name)) "&")) 
%  presentation (shell-command (concat "xpdf -fullscreen " (replace-regexp-in-string "tex$" "pdf"(buffer-file-name)) "&"))
%  presentation (shell-command (concat "evince " (replace-regexp-in-string "tex$" "pdf"(buffer-file-name)) "&"))

%http://www.naney.org/diki/dk/hyperref.html
%日本語EUC系環境の時
\AtBeginDvi{\special{pdf:tounicode EUC-UCS2}}
%シフトJIS系環境の時
%\AtBeginDvi{\special{pdf:tounicode 90ms-RKSJ-UCS2}}

\newenvironment{commandlinesmall}%
{\VerbatimEnvironment
  \begin{Sbox}\begin{minipage}{1.0\hsize}\begin{fontsize}{8}{8} \begin{BVerbatim}}%
{\end{BVerbatim}\end{fontsize}\end{minipage}\end{Sbox}
  \setlength{\fboxsep}{8pt}
% start on a new paragraph

\vspace{6pt}% skip before
\fcolorbox{dancerdarkblue}{dancerlightblue}{\TheSbox}

\vspace{6pt}% skip after
}
%end of commandlinesmall

\title{東京エリアDebian勉強会}
\subtitle{第144回 2016年10月度}
\author{杉本 典充}
\date{2016年10月15日}
\logo{\includegraphics[width=8cm]{image200607/openlogo-light.eps}}

\begin{document}

\begin{frame}
\titlepage{}
\end{frame}

\begin{frame}{Agenda}
 \begin{minipage}[t]{0.45\hsize}
  \begin{itemize}
  \item 最近あったDebian関連のイベント報告
	\begin{itemize}
	\item 第143回東京エリアDebian勉強会
	\end{itemize}
  \item 事前課題発表
  \end{itemize}
 \end{minipage}
 \begin{minipage}[t]{0.45\hsize}
  \begin{itemize}
   \item Debian Trivia Quiz
  \end{itemize}
 \end{minipage}
\end{frame}

\section{イベント報告}
\emtext{イベント報告}

\begin{frame}{第143回東京エリアDebian勉強会}

\begin{itemize}
\item 場所は朝日ネットさんをお借りして開催しました
\item 参加者は6名
\item セミナ内容は、以下内容を発表しました
  \begin{enumerate}
  \item iwamatsuさん「DEP5 Machine-readable debian/copyright 再考」
  \end{enumerate}
\item 残りの時間でhack timeを行い、成果発表をしました
\item \url{https://tokyodebianmeeting.titanpad.com/ep/pad/export/143/latest?format=txt}
\end{itemize} 
\end{frame}

\section{事前課題}
\emtext{事前課題}
{\footnotesize
 \begin{prework}{ nabaua }
  \begin{enumerate}
  \item 賑やかしで～
  \end{enumerate}
\end{prework}

\begin{prework}{ dictoss }
  \begin{enumerate}
  \item tilegx環境をつくるべくLFSする
  \end{enumerate}
\end{prework}

\begin{prework}{ kamagasako }
  \begin{enumerate}
  \item preseedのおさらい、Vagrantと組み合わせてVirtualBox上にdebianを自動インストール
  \end{enumerate}
\end{prework}

\begin{prework}{ Roger Shimizu }
  \begin{enumerate}
  \item パッケージメインテイン
  \end{enumerate}
\end{prework}

\begin{prework}{ koedoyoshida }
  \begin{enumerate}
  \item 所用で途中で抜けます。
  \end{enumerate}
\end{prework}

}

\subsection{問題}

%; whizzy-master ../debianmeetingresume201311.tex
% $B0J>e$N@_Dj$r$7$F$$$k$?$a!"$3$N%U%!%$%k$G(B M-x whizzytex $B$9$k$H!"(Bwhizzytex$B$,MxMQ$G$-$^$9!#(B
%

\santaku
{HP$B%(%s%?!<%W%i%$%:<R$+$i(BDebian Project$B$X%5!<%P$r4sIU$7$F$$$@$-$^$7$?!#$3$N%5!<%P$r;H$C$F(Bdebian.org$B$N%5!<%S%9$N:\$;BX$($,<B;\$5$l$^$9!#$5$F!"$I$NAH$_9g$o$;$,@5$7$$$G$7$g$&$+(B}
{ftp-master.debian.org$B$H(Bwww.debian.org}
{ftp-master.debian.org$B$H(Bsecurity.debian.org}
{www.debian.org$B$H(Bsecurity.debian.org}
{B}
{$B4sIU$$$?$@$$$?%5!<%P$O%+%J%@!"%"%a%j%+!"%*!<%9%H%i%j%"$K$"$k%5!<%P$N99?7$KMxMQ$9$kM=Dj$G$9!#6qBNE*$K$O(BDebian$B$N<gMW%5!<%S%9$G$"$k(Bftp-master.debian.org$B$H(Bsecurity.debian.org$B$K3d$jEv$F$7!"$=$NB>$N%3%"%5!<%S%9$N8~>e$K$bLrN)$F$kM=Dj$G$9!#(B\url{https://www.debian.org/News/2016/20161003}}


\section{セミナー1}
\emtext{セミナー1}

\section{Hack time}
\emtext{Hack time}

\begin{frame}{成果を記入下さい!}

  今回、Hack time時の成果を記録に残してみます。

\url{https://tokyodebianmeeting.titanpad.com/144}

に皆さんアクセス頂き（認証不要です）、各自成果を17:30までに記録するようにお願いします。

\end{frame}
  
\section{今後のイベント}
\emtext{今後のイベント}
\begin{frame}{今後のイベント}
  \begin{itemize}
  \item 11/5(土),6(日) OSC 2016 Tokyo/Fall
    \begin{itemize}  
    \item セミナー、展示は11/5
    \end{itemize}
  \item 11/19(土) 東京エリアDebian勉強会(予定)
    \begin{itemize}
    \item 発表者募集中
    \end{itemize}
  \item 12/10(土) Mini Debian Conference Japan 2016
    \begin{itemize}
    \item \url{http://miniconf.debian.or.jp/}
    \item 発表者募集中
    \end{itemize}    
  \item 12月の東京エリアDebian勉強会はお休みの予定です
  \item 2017/1/21(土) 東京エリアDebian勉強会(予定)
  \end{itemize}
\end{frame}

\end{document}

;;; Local Variables: ***
;;; outline-regexp: "\\([ 	]*\\\\\\(documentstyle\\|documentclass\\|emtext\\|section\\|begin{frame}\\)\\*?[ 	]*[[{]\\|[]+\\)" ***
;;; End: ***
